\chapter{Expérimentations et résultats}
\label{chap:experimentations}

\section{Introduction}\label{sec:c6-intro}

Ce chapitre réunit les expérimentations du mémoire. Il évalue d'abord l'extraction multimodale (chapitre~\ref{chap:extraction}) : choix du modèle, performance de référence et nature des erreurs. Il évalue ensuite la chaîne TRACE (chapitre~\ref{chap:architecture}) appliquée à cette extraction : apport de chaque étage, mode standard et mode précision. Toutes les mesures portent sur les mêmes dossiers, avec le même Gold Standard et la même validation croisée.

\section{Protocole expérimental}\label{sec:c6-protocole}

Les évaluations portent sur les 49 dossiers du Gold Standard final (291 pages, 12\,396 entités, 5\,546 relations). L'extraction suit le pipeline du chapitre~\ref{chap:extraction}, à température~0. Tout ce qui s'appuie sur la référence est appris en \textbf{validation croisée à deux plis} : les dossiers sont triés par ordre alphabétique de leur identifiant, puis répartis un sur deux (25 et 24 dossiers), et chaque moitié est traitée avec ce qui a été appris sur l'autre (section~\ref{sec:c5-validation}). Les performances rapportées valent donc pour des dossiers non vus.

Les extractions sont comparées au Gold Standard avec les trois modes d'appariement de la section~\ref{sec:c4-matching}, et les mesures de la section~\ref{sec:c4-metriques} : précision, rappel, F1-score et taux d'hallucination $H = 1 - \text{Précision}$. Pour TRACE, ces mesures sont calculées à la sortie de chaque étage. Sauf mention contraire, les résultats sont donnés en mode Relaxed et en micro-moyenne, en pourcentage. Le mode Relaxed juge le contenu d'un élément plutôt que sa forme exacte : un modèle génératif reformule souvent légèrement une mention juste (\og noradrénaline IVSE \fg{} pour \og noradrénaline \fg{}), et une correspondance exacte la compterait à tort comme fausse~\cite{wadhwa2023revisiting}. Il reste exigeant, puisque l'élément doit se trouver sur la même page et avoir le même type que dans la référence. La micro-moyenne donne le même poids à chaque élément et mesure donc la qualité du graphe sur l'ensemble du corpus.

\section{Choix du modèle d'extraction}\label{sec:c6-choix}

Les deux modèles ont été comparés avec la même consigne, à température~0, sur 43 dossiers annotés avant la finalisation du guideline. Cette étude a été présentée à la conférence JDSE~2026~\cite{dridi2026jdse} (tableau~\ref{tab:c6-preliminaire}).

\begin{table}[H]
\centering
\caption{Comparaison des deux modèles (43 dossiers, Gold Standard préliminaire, mode Relaxed, en \%)}\label{tab:c6-preliminaire}
\small
\begin{tabular}{llccc}
\toprule
& \textbf{Modèle} & \textbf{Précision} & \textbf{Rappel} & \textbf{F1} \\
\midrule
Entités & Qwen2.5-VL-7B-Instruct & 64,10 & 63,67 & 63,88 \\
 & Mistral Small 4 & \textbf{67,38} & \textbf{71,89} & \textbf{69,57} \\
\midrule
Relations & Qwen2.5-VL-7B-Instruct & 50,83 & 30,31 & 37,98 \\
 & Mistral Small 4 & \textbf{67,42} & \textbf{58,17} & \textbf{62,43} \\
\bottomrule
\end{tabular}
\end{table}

Mistral Small~4 dépasse Qwen2.5-VL sur toutes les mesures, modérément sur les entités (+5,7~points de F1) et très nettement sur les relations (+24,5~points). L'écart vient surtout du rappel des relations, presque deux fois plus élevé : Qwen2.5-VL retrouve des entités, mais en relie beaucoup moins. Le critère principal, fixé avant l'expérience, est la qualité des relations, puisque dans un graphe de connaissances la connaissance est portée par les liens~\cite{hogan2021knowledge} : savoir qu'un patient a reçu de la noradrénaline n'a de valeur que si l'on sait à quelle dose et pour quelle défaillance. \textbf{Mistral Small~4 a donc été retenu}.

Ce choix a une contrepartie : Mistral Small~4 a été utilisé par une API externe pour la version finale de l'extraction. Les précautions prises ne suffisent pas à établir la conformité de ce traitement au RGPD (section~\ref{sec:c4-donnees}). Ses poids étant publiés, il peut être exécuté localement, avec des performances qui devraient être comparables.

\section{Performance de référence}\label{sec:c6-reference}

Le tableau~\ref{tab:c6-reference} donne les performances de Mistral Small~4 à la sortie du pipeline final, avant TRACE. C'est le point de départ de toutes les mesures qui suivent.

\begin{table}[H]
\centering
\caption{Performance de référence avant TRACE (49 dossiers, validation croisée, mode Relaxed, micro-moyenne, en \%)}\label{tab:c6-reference}
\small
\begin{tabular}{lcccc}
\toprule
& \textbf{Précision} & \textbf{Rappel} & \textbf{F1} & \textbf{H} \\
\midrule
Entités & 66,20 & 68,45 & 67,31 & 33,8 \\
Relations & 54,77 & 45,46 & 49,68 & 45,2 \\
\bottomrule
\end{tabular}
\end{table}

Un tiers des entités extraites et près de la moitié des relations ne sont pas justifiées par la référence. Les relations sont plus difficiles, ce qui est cohérent avec la littérature~\cite{wadhwa2023revisiting} : une relation suppose deux entités justes, correctement reliées. Ces résultats sont plus faibles que ceux de l'étude préliminaire, car la référence finale est plus exigeante (revue par le médecin, complétée par les absences) et toutes les composantes apprises sont désormais évaluées sur des dossiers non vus.

\paragraph{Analyse.} Le passage du mode Strict au mode Boundary-Relaxed fait gagner environ 14~points de F1 sur les entités et 20 sur les relations : une grande partie des écarts sont des différences de délimitation, non de contenu. Le passage au mode Relaxed n'ajoute ensuite qu'environ un point, et les micro- et macro-moyennes diffèrent de moins d'un point : les conclusions ne dépendent ni du mode ni de quelques dossiers. Par type, les entités à forme régulière sont les mieux extraites (biomarqueurs, micro-organismes, traitements), et les entités d'interprétation clinique les moins bien (diagnostics, évolution, foyers infectieux). Les relations les mieux extraites sont aussi les plus explicites : l'administration d'un traitement au patient et le lien entre un traitement et sa posologie.

\section{Nature des erreurs}\label{sec:c6-erreurs}\label{sec:c6-discussion}

Les erreurs relèvent de trois familles. Les \textbf{éléments en trop} définissent le taux d'hallucination : entités inventées ou dont la valeur n'est pas écrite, mentions mal délimitées, éléments de contexte non annotés (un symptôme des antécédents, un traitement envisagé) et relations déduites de la seule présence de deux entités. Les \textbf{erreurs de type} portent sur des catégories voisines (signe vital et symptôme, diagnostic et foyer) : ce sont des hallucinations terminologiques. Les \textbf{oublis}, enfin, se concentrent dans les tableaux de biologie et les ordonnances ; ils ne sont pas des hallucinations, mais limitent ce qu'un post-traitement peut améliorer.

Cette évaluation fixe donc l'objectif de TRACE : la marge porte d'abord sur la \textbf{précision}, c'est-à-dire sur les hallucinations. Le rappel, lui, est borné par l'extraction, sauf dans les blocs à forme fixe, où la récupération structurée peut recopier les lignes oubliées.

\section{Apport de chaque étage de TRACE}\label{sec:c6-trace-etages}

Le tableau~\ref{tab:c6-etages} donne les performances à la sortie de chaque étage, en mode standard. Les étages~3, 5 et~6 annotent sans retirer : leurs mesures sont pratiquement identiques à celles de l'étage précédent, et leur apport se mesure à l'étage~7, qui utilise leurs annotations.

\begin{table}[H]
\centering
\caption{Performances à la sortie des étages de TRACE (49 dossiers, mode Relaxed, micro-moyenne, en \%)}\label{tab:c6-etages}
\small
\begin{tabular}{lcccccccc}
\toprule
& \multicolumn{4}{c}{\textbf{Entités}} & \multicolumn{4}{c}{\textbf{Relations}} \\
\cmidrule(lr){2-5}\cmidrule(lr){6-9}
\textbf{Sortie} & \textbf{P} & \textbf{R} & \textbf{F1} & \textbf{H} & \textbf{P} & \textbf{R} & \textbf{F1} & \textbf{H} \\
\midrule
Extraction (avant TRACE) & 66,20 & 68,45 & 67,31 & 33,8 & 54,77 & 45,46 & 49,68 & 45,2 \\
Étage 2 · ancrage & 66,20 & 68,45 & 67,31 & 33,8 & 55,76 & 45,36 & 50,02 & 44,2 \\
Étage 4 · SGCE & 66,20 & 68,46 & 67,31 & 33,8 & 55,75 & 45,38 & 50,03 & 44,2 \\
Étage 4b · récupération & 65,92 & 68,88 & 67,37 & 34,1 & 55,73 & 46,75 & 50,85 & 44,3 \\
Étage 7 · arbitrage standard & 66,32 & 67,96 & 67,13 & 33,7 & 56,31 & 45,56 & 50,37 & 43,7 \\
\bottomrule
\end{tabular}
\end{table}

En reconstituant les nombres d'éléments à partir du rappel et de la précision~\cite{sokolova2009}, on peut juger chaque étage sur ce qu'il retire ou ajoute. L'\textbf{ancrage} retire environ 92 relations, dont environ 6 seulement étaient justes : plus de neuf sur dix étaient des hallucinations. \textbf{SGCE} modifie peu les mesures, car il répare la forme du graphe, que le F1 calculé élément par élément voit peu ; surtout, il ne dégrade aucune mesure. La \textbf{récupération} ajoute environ 138 relations, dont 76 justes, et apporte le plus fort gain de F1 de la chaîne (+0,8~point sur les relations). L'\textbf{arbitrage} retire environ 250 entités et 164 relations ; le rappel baisse d'environ un point, surtout à cause de la règle E1, qui retire des reformulations parfois justes mais intraçables.

Au total, le mode standard réduit le taux d'hallucination des relations de 45,2 à 43,7~\% et améliore leur F1, en laissant les entités pratiquement inchangées. Ce gain est modeste par construction : le mode standard n'apprend rien de la référence et ne retire que ce que l'on peut démontrer faux ou intraçable. C'est un plancher sûr.

\section{Mode précision}\label{sec:c6-precision}

Le tableau~\ref{tab:c6-precision} donne les performances du mode précision en validation croisée, selon la part $k$ d'éléments conservés.

\begin{table}[H]
\centering
\caption{Mode précision de TRACE en validation croisée, selon la part $k$ d'éléments conservés (en \%)}\label{tab:c6-precision}
\small
\begin{tabular}{lcccccccc}
\toprule
& \multicolumn{4}{c}{\textbf{Entités}} & \multicolumn{4}{c}{\textbf{Relations}} \\
\cmidrule(lr){2-5}\cmidrule(lr){6-9}
\textbf{$k$} & \textbf{P} & \textbf{R} & \textbf{F1} & \textbf{H} & \textbf{P} & \textbf{R} & \textbf{F1} & \textbf{H} \\
\midrule
\emph{Avant TRACE} & \emph{66,20} & \emph{68,45} & \emph{67,31} & \emph{33,8} & \emph{54,77} & \emph{45,46} & \emph{49,68} & \emph{45,2} \\
95~\% & 68,01 & 67,22 & 67,62 & 32,0 & 57,70 & 45,15 & \textbf{50,66} & 42,3 \\
90~\% & 69,83 & 66,30 & 68,02 & 30,2 & 58,98 & 44,27 & 50,58 & 41,0 \\
80~\% & 74,10 & 62,92 & \textbf{68,05} & 25,9 & 62,19 & 41,47 & 49,76 & 37,8 \\
70~\% & 78,15 & 58,03 & 66,60 & 21,9 & 65,49 & 37,57 & 47,75 & 34,5 \\
50~\% & 86,98 & 44,51 & 58,89 & 13,0 & 72,98 & 21,82 & 33,60 & 27,0 \\
30~\% & \textbf{91,60} & 27,81 & 42,66 & \textbf{8,4} & \textbf{79,11} & 6,20 & 11,49 & \textbf{20,9} \\
\bottomrule
\end{tabular}
\end{table}

Plus $k$ diminue, plus la précision augmente et le taux d'hallucination baisse, de façon régulière : le score appris ordonne correctement les éléments~\cite{chow1970reject,geifman2017selective}. \textbf{De 95 à 80~\%}, le vote réduit les hallucinations sans perte de qualité globale : à $k = 80~\%$, le taux d'hallucination passe de 33,8 à 25,9~\% sur les entités et de 45,2 à 37,8~\% sur les relations, à F1 égal ou supérieur. En nombre d'éléments, les entités fausses diminuent de 37~\% pour seulement 8~\% d'entités justes perdues : le vote retire plus de deux éléments faux pour un élément juste. \textbf{À 30~\%}, le graphe devient très fiable (91,6~\% de précision sur les entités), mais très incomplet ; ce réglage ne convient qu'à l'extraction d'un petit noyau d'informations sûres. Le réglage se choisit donc selon l'usage : exhaustivité pour explorer un dossier, fiabilité pour l'aide à la décision.

\section{Discussion}\label{sec:c6-bilan}\label{sec:c6-discussion-generale}

TRACE réduit les hallucinations sans dégrader le graphe lorsque le réglage est modéré, donne le choix du niveau d'exigence et garde chaque décision traçable. Ses gains de F1 restent modestes, et ce n'est pas un défaut de conception : un post-traitement ne peut pas retrouver ce que le modèle n'a jamais extrait, si bien que sa marge porte avant tout sur la précision. Tous les résultats ont été obtenus en validation croisée : un vote appris et évalué sur les mêmes dossiers donnerait des chiffres plus élevés, mais trompeurs~\cite{hastie2009elements}. Enfin, le taux d'hallucination majore la part d'informations réellement inventées : une partie des faux positifs sont des éléments écrits dans le dossier mais non annotés, ou des différences de délimitation.

\section{Conclusion}\label{sec:c6-conclusion}

L'évaluation de l'extraction a conduit à retenir Mistral Small~4, qui atteint en validation croisée un F1 de 67,3~\% sur les entités et de 49,7~\% sur les relations, avec des taux d'hallucination de 33,8~\% et de 45,2~\%. L'évaluation de TRACE montre que chaque étage joue le rôle prévu sans dégrader le résultat, que le mode standard réduit les hallucinations des relations en préservant le reste du graphe, et que le mode précision offre un compromis réglable : à $k = 80~\%$, le taux d'hallucination passe de 33,8 à 25,9~\% sur les entités et de 45,2 à 37,8~\% sur les relations, à F1 égal ou supérieur.
